% 分野別類題: リカッチ方程式（Riccati equation）
% タイトルのみ英語、本文は日本語。特解 y1 を与え、y = y1 + 1/u でベルヌーイ／線形に落とす定番 3 問。
% コンパイル: TEXINPUTS="../../../00_common:$TEXINPUTS" latexmk -xelatex problems.tex
\documentclass[a4paper,11pt]{article}
\usepackage{preamble}

\begin{document}

\kamokumei{類題演習 --- Riccati Differential Equations}

\shijibun{次の常微分方程式を解き、導出過程を示せ。（ヒント: リカッチ方程式は特解 $y_{1}$ が1つ分かれば、$y = y_{1} + \dfrac{1}{u}$ の置換により $u$ についての1階線形方程式に変換できる。）}

\shomon{問1.}{$y_{1} = x$ が次の微分方程式の特解であることを確かめ、一般解を求めよ。}
\[
\frac{dy}{dx} = 1 + x^{2} - 2xy + y^{2}
\]

\shomon{問2.}{$y_{1} = \dfrac{1}{x}$ が次の微分方程式の特解であることを確かめ、一般解を求めよ。}
\[
\frac{dy}{dx} = -\frac{1}{x^{2}} - \frac{y}{x} + y^{2}
\]

\shomon{問3.}{$y_{1} = e^{x}$ が次の微分方程式の特解であることを用いて、一般解を求めよ。}
\[
\frac{dy}{dx} = e^{-x}y^{2} + y - e^{x}
\]

\end{document}
